\begin{exercise}
\label{ex:normalization-curry-howard-01}
Type the term \(\lambda x{:}A.\lambda y{:}B.x\), and give the corresponding
natural-deduction derivation.  Identify the structural weakening implicit in
the unused variable \(y\).
\end{exercise}

\begin{exercise}
\label{ex:normalization-curry-howard-02}
Prove subject reduction for a single outermost beta redex.  State the typing
substitution lemma exactly and mark where capture avoidance is used.
\end{exercise}

\begin{definition}[Abstract reduction relation and normalization]
An \emph{abstract reduction relation} is a binary relation \(\to\) on a
set.  An element is \emph{strongly normalizing} iff no infinite
\(\to\)-sequence starts from it; the relation is strongly normalizing iff
every element is.  An element is \emph{weakly normalizing} iff some finite
\(\to\)-sequence from it ends in an element with no \(\to\)-successor;
the relation is weakly normalizing iff every element is.  Beta reduction on
well-typed terms is an abstract reduction relation, and the
strong-normalization theorem above states that it is strongly normalizing.
\end{definition}

\begin{exercise}
\label{ex:normalization-curry-howard-03}
Explain why a strongly normalizing abstract reduction relation is weakly
normalizing, and exhibit an abstract reduction relation that is weakly but
not strongly normalizing.
\end{exercise}

